% GENERATED FILE. Edit canonical lessons, then run npm run generate:books.
% canonical-source-hash: fnv1a64:939f6ce776a5355f
% canonical-lessons: TA-C312-virivaka, TA-C312-neratiyaka, TA-C312-varni, TA-C312-oppi, TA-C312-vaymoli

\chapter{Telling It Clearly}
\label{ch:ta-telling-it-clearly}
\hlchaptermodality{312}

\begin{chapteropening}
\textbf{By the end of this chapter:} \emph{I can tell a story in detail and in my own words.}

Complete the last lesson of chapter 312: I can tell a story in detail and in my own words.
\end{chapteropening}

\section[\texorpdfstring{\ta{விரிவாக} (virivāka)}{விரிவாக (virivāka)}]{\ta{விரிவாக} (virivāka) — in detail}

\label{lesson:TA-C312-virivaka}



\begin{quote}
Before the new one: say the Tamil for an uprising, then the Tamil for a proclamation.
\end{quote}

\subsection*{You'll want to know: \ta{விரிவாக}}
\textbf{\ta{விரிவாக}} — \emph{virivāka} — \textquotedblleft{}in detail\textquotedblright{}.

\begin{quote}
\ta{சாட்சி} \ta{நடந்ததை} \ta{விரிவாக} \ta{விளக்கினார்}.
\end{quote}

\begin{quote}
\emph{cāṭci naṭantatai virivāka viḷakkiṉār.}
\end{quote}

\begin{quote}
The witness explained what happened in detail.
\end{quote}

-\ta{ஆக} makes the adverb; the opposite is \ta{சுருக்கமாக}, briefly.

\begin{cousinweb}
\textbf{\ta{விரிவாக}} (\emph{virivāka}) is built on native \ta{விரிவு}, \textquotedblleft{}spread, expansion\textquotedblright{}, from \ta{விரி}, \textquotedblleft{}to expand, spread out\textquotedblright{} (DEDR 5411), with the adverb ending -\ta{ஆக}.
\end{cousinweb}

\begin{grammarlens}[title={What you've built}]
One more word for telling what happened.
\end{grammarlens}

\subsection*{Your turn}
\begin{itemize}
\raggedright
  \item \emph{Say it:} \emph{virivāka}
  \item \emph{Say it:} \emph{virivāka}, once more
  \item \emph{Say it:} \emph{cāṭci naṭantatai virivāka viḷakkiṉār.} — The witness explained what happened in detail.
  \item \emph{Recall:} say the Tamil for an uprising, then the Tamil for a proclamation, then say \emph{virivāka} again
\end{itemize}

\begin{tcolorbox}[breakable,colback=teal!4,colframe=teal!35!black,title={Before you move on}]
What does \ta{விரிவாக} mean? (\textquotedblleft{}In detail\textquotedblright{}.) Say it once more.
\end{tcolorbox}
\section[\texorpdfstring{\ta{நேரடியாக} (nēraṭiyāka)}{நேரடியாக (nēraṭiyāka)}]{\ta{நேரடியாக} (nēraṭiyāka) — directly; in person}

\label{lesson:TA-C312-neratiyaka}



\begin{quote}
Before the new one: say the Tamil for a proclamation, then the Tamil for in detail.
\end{quote}

\subsection*{You'll want to know: \ta{நேரடியாக}}
\textbf{\ta{நேரடியாக}} — \emph{nēraṭiyāka} — \textquotedblleft{}directly; in person\textquotedblright{}.

\begin{quote}
\ta{அவர்} \ta{என்னிடம்} \ta{நேரடியாகக்} \ta{கேட்டார்}.
\end{quote}

\begin{quote}
\emph{avar eṉṉiṭam nēraṭiyākak kēṭṭār.}
\end{quote}

\begin{quote}
He asked me directly.
\end{quote}

\ta{நேரடியாக} \ta{ஒளிபரப்பு}, to broadcast live.

\begin{cousinweb}
\textbf{\ta{நேரடியாக}} (\emph{nēraṭiyāka}) is built on native \ta{நேர்}, \textquotedblleft{}straightness, directness\textquotedblright{} (DEDR 3772), with \ta{அடி} and the adverb ending -\ta{ஆக}.
\end{cousinweb}

\begin{grammarlens}[title={What you've built}]
One more word for telling what happened.
\end{grammarlens}

\subsection*{Your turn}
\begin{itemize}
\raggedright
  \item \emph{Say it:} \emph{nēraṭiyāka}
  \item \emph{Say it:} \emph{nēraṭiyāka}, once more
  \item \emph{Say it:} \emph{avar eṉṉiṭam nēraṭiyākak kēṭṭār.} — He asked me directly.
  \item \emph{Recall:} say the Tamil for a proclamation, then the Tamil for in detail, then say \emph{nēraṭiyāka} again
\end{itemize}

\begin{tcolorbox}[breakable,colback=teal!4,colframe=teal!35!black,title={Before you move on}]
What does \ta{நேரடியாக} mean? (\textquotedblleft{}Directly; in person\textquotedblright{}.) Say it once more.
\end{tcolorbox}
\section[\texorpdfstring{\ta{வர்ணி} (varṇi)}{வர்ணி (varṇi)}]{\ta{வர்ணி} (varṇi) — to describe vividly}

\label{lesson:TA-C312-varni}



\begin{quote}
Before the new one: say the Tamil for in detail, then the Tamil for directly.
\end{quote}

\subsection*{You'll want to know: \ta{வர்ணி}}
\textbf{\ta{வர்ணி}} — \emph{varṇi} — \textquotedblleft{}to describe vividly\textquotedblright{}.

\begin{quote}
\ta{தாத்தா} \ta{அந்தப்} \ta{பழைய} \ta{திருவிழாவை} \ta{அழகாக} \ta{வர்ணித்தார்}.
\end{quote}

\begin{quote}
\emph{tāttā antap paḻaiya tiruviḻāvai aḻakāka varṇittār.}
\end{quote}

\begin{quote}
Grandfather described that old festival beautifully.
\end{quote}

Past (I …): \ta{வர்ணித்தேன்}, \emph{varṇittēṉ}.

The noun \ta{வர்ணனை} is a description or a sports commentary.

\begin{cousinweb}
\textbf{\ta{வர்ணி}} (\emph{varṇi}) is a Tamil verb built on Sanskrit \emph{varṇana}, \textquotedblleft{}painting, description\textquotedblright{} (Monier-Williams).
\end{cousinweb}

\begin{grammarlens}[title={What you've built}]
One more word for telling what happened.
\end{grammarlens}

\subsection*{Your turn}
\begin{itemize}
\raggedright
  \item \emph{Say it:} \emph{varṇi}
  \item \emph{Say it:} \emph{varṇi}, once more
  \item \emph{Say it:} \emph{tāttā antap paḻaiya tiruviḻāvai aḻakāka varṇittār.} — Grandfather described that old festival beautifully.
  \item \emph{Recall:} say the Tamil for in detail, then the Tamil for directly, then say \emph{varṇi} again
\end{itemize}

\begin{tcolorbox}[breakable,colback=teal!4,colframe=teal!35!black,title={Before you move on}]
What does \ta{வர்ணி} mean? (\textquotedblleft{}To describe vividly\textquotedblright{}.) Say it once more.
\end{tcolorbox}
\section[\texorpdfstring{\ta{ஒப்பி} (oppi)}{ஒப்பி (oppi)}]{\ta{ஒப்பி} (oppi) — to recite}

\label{lesson:TA-C312-oppi}



\begin{quote}
Before the new one: say the Tamil for directly, then the Tamil for to describe vividly.
\end{quote}

\subsection*{You'll want to know: \ta{ஒப்பி}}
\textbf{\ta{ஒப்பி}} — \emph{oppi} — \textquotedblleft{}to recite\textquotedblright{}.

\begin{quote}
\ta{மாணவன்} \ta{பாடத்தை} \ta{மனப்பாடமாக} \ta{ஒப்பித்தான்}.
\end{quote}

\begin{quote}
\emph{māṇavaṉ pāṭattai maṉappāṭamāka oppittāṉ.}
\end{quote}

\begin{quote}
The student recited the lesson by heart.
\end{quote}

Past (I …): \ta{ஒப்பித்தேன்}, \emph{oppittēṉ}.

Said of reciting a lesson or a poem from memory.

\begin{cousinweb}
\textbf{\ta{ஒப்பி}} (\emph{oppi}), \textquotedblleft{}to recite, to hand over\textquotedblright{}, is a Tamil verb whose origin is not settled in the sources at hand.
\end{cousinweb}

\begin{grammarlens}[title={What you've built}]
One more word for telling what happened.
\end{grammarlens}

\subsection*{Your turn}
\begin{itemize}
\raggedright
  \item \emph{Say it:} \emph{oppi}
  \item \emph{Say it:} \emph{oppi}, once more
  \item \emph{Say it:} \emph{māṇavaṉ pāṭattai maṉappāṭamāka oppittāṉ.} — The student recited the lesson by heart.
  \item \emph{Recall:} say the Tamil for directly, then the Tamil for to describe vividly, then say \emph{oppi} again
\end{itemize}

\begin{tcolorbox}[breakable,colback=teal!4,colframe=teal!35!black,title={Before you move on}]
What does \ta{ஒப்பி} mean? (\textquotedblleft{}To recite\textquotedblright{}.) Say it once more.
\end{tcolorbox}
\section[\texorpdfstring{\ta{வாய்மொழி} (vāymoḻi)}{வாய்மொழி (vāymoḻi)}]{\ta{வாய்மொழி} (vāymoḻi) — oral tradition; word of mouth}

\label{lesson:TA-C312-vaymoli}



\begin{quote}
Before the new one: say the Tamil for to describe vividly, then the Tamil for to recite.
\end{quote}

\subsection*{You'll want to know: \ta{வாய்மொழி}}
\textbf{\ta{வாய்மொழி}} — \emph{vāymoḻi} — \textquotedblleft{}oral tradition; word of mouth\textquotedblright{}.

\begin{quote}
\ta{இந்தக்} \ta{கதைகள்} \ta{வாய்மொழியாக} \ta{வந்தவை}.
\end{quote}

\begin{quote}
\emph{intak kataikaḷ vāymoḻiyāka vantavai.}
\end{quote}

\begin{quote}
These stories have come down by word of mouth.
\end{quote}

\ta{வாய்மொழியாக}, orally; \ta{வாய்}, mouth + \ta{மொழி}, word.

\begin{cousinweb}
\textbf{\ta{வாய்மொழி}} (\emph{vāymoḻi}) is a Tamil compound built from \ta{வாய்}, \textquotedblleft{}mouth\textquotedblright{} (DEDR 5352), and \ta{மொழி}, \textquotedblleft{}word, saying\textquotedblright{} (DEDR 4989).
\end{cousinweb}

\begin{grammarlens}[title={What you've built}]
One more word for telling what happened.
\end{grammarlens}

\subsection*{Your turn}
\begin{itemize}
\raggedright
  \item \emph{Say it:} \emph{vāymoḻi}
  \item \emph{Say it:} \emph{vāymoḻi}, once more
  \item \emph{Say it:} \emph{intak kataikaḷ vāymoḻiyāka vantavai.} — These stories have come down by word of mouth.
  \item \emph{Recall:} say the Tamil for to describe vividly, then the Tamil for to recite, then say \emph{vāymoḻi} again
\end{itemize}

\begin{tcolorbox}[breakable,colback=teal!4,colframe=teal!35!black,title={Before you move on}]
What does \ta{வாய்மொழி} mean? (\textquotedblleft{}Oral tradition; word of mouth\textquotedblright{}.) Say it once more.
\end{tcolorbox}
